\section{Organic Chemistry, Functional Groups, Chirality, and Polymers}

\subsection{Ctrl+F keywords: functional group, alcohol, ether, ester, amide, carboxylic acid, aromatic, chirality, isomer, polymerization}
Use for identifying organic functional groups, aromatic rings, isomers, chiral centers, amino acids, peptides, and polymer repeat units.

\subsection{Symbols and notation}
\referencetable{tab:organic:symbols}{Symbols and notation for organic chemistry.}
\begin{tabular}{@{}>{$}l<{$} p{10.8cm}@{}}
\toprule
\text{Symbol} & \text{Meaning and usage}\\
\midrule
n & Number of carbon atoms in general hydrocarbon formulas; polymerization degree when used with polymer masses.\\
\mathrm C_n\mathrm H_{2n+2} & General formula for an acyclic alkane.\\
\mathrm C_n\mathrm H_{2n},\ \mathrm C_n\mathrm H_{2n-2} & General formulas for acyclic alkenes and alkynes.\\
\mathrm{DBE} & Degree of unsaturation: total rings and \(\pi\)-bonds.\\
C,\ H,\ N,\ X & Atom counts in the DBE formula; \(X\) counts halogens.\\
\mathrm R,\ \mathrm{R'} & Generic organic groups; primes distinguish groups that need not be identical.\\
\mathrm{C^*} & A tetrahedral stereogenic carbon attached to four different groups.\\
M_{\mathrm{polymer}},\ M_{\mathrm{repeat}} & Polymer molar mass and repeat-unit molar mass.\\
\bottomrule
\end{tabular}

\subsection{Hydrocarbons and unsaturation}

\subsubsection{Alkane, alkene, alkyne formulas}
\begin{equation}
\label{eq:organic:hydrocarbon-formulas}
\text{alkane: }\mathrm{C}_n\mathrm{H}_{2n+2}, \qquad
\text{alkene: }\mathrm{C}_n\mathrm{H}_{2n}, \qquad
\text{alkyne: }\mathrm{C}_n\mathrm{H}_{2n-2}
\end{equation}
Alkenes contain \(\ce{C=C}\) and alkynes contain \(\ce{C#C}\). Rings reduce hydrogen count like one double bond.

\subsubsection{Degree of unsaturation, double bond equivalent}
\begin{equation}
\label{eq:organic:dbe}
\mathrm{DBE}=\frac{2C+2+N-H-X}{2}
\end{equation}
Oxygen does not affect DBE. Each ring or \(\pi\)-bond counts as one DBE; benzene has four DBE.

\subsubsection{Aromatic ring recognition}
\begin{equation}
\label{eq:organic:aromatic-motif}
\ce{C6H6}\ \text{ring motif}, \qquad \text{phenyl group: }\mathrm{C_6H_5{-}}
\end{equation}
A benzene ring is planar and aromatic. Aromatic rings often occur with alcohol, carboxylic acid, amide, or ester groups.

\subsection{Functional group lookup}

\subsubsection{Functional group recognition table}
\referencetable{tab:organic:functional-groups}{Functional groups recognized from structural motifs.}
\begin{tabular}{@{}ll@{}}
\toprule
\textbf{Group} & \textbf{Structural motif}\\
\midrule
alcohol & \(\mathrm{R{-}OH}\)\\
ether & \(\mathrm{R{-}O{-}R'}\)\\
aldehyde & \(\mathrm{R{-}CHO}\)\\
ketone & \(\mathrm{R{-}C(=O){-}R'}\)\\
carboxylic acid & \(\mathrm{R{-}COOH}\)\\
ester & \(\mathrm{R{-}C(=O){-}O{-}R'}\)\\
amine & \(\mathrm{R{-}NH_2},\ \mathrm{R_2NH},\ \mathrm{R_3N}\)\\
amide & \(\mathrm{R{-}C(=O){-}NH_2},\ \mathrm{R{-}C(=O){-}NR_2}\)\\
nitrile & \(\mathrm{R{-}C{\equiv}N}\)\\
\bottomrule
\end{tabular}
Carbonyl-containing groups are distinguished by what is attached directly to the carbonyl carbon.

\subsubsection{Ester, ether, amide, amine distinction}
\begin{equation}
\label{eq:organic:heteroatom-distinction}
\text{ester: }\mathrm{C(=O)O}, \qquad
\text{ether: }\mathrm{C{-}O{-}C}, \qquad
\text{amide: }\mathrm{C(=O)N}, \qquad
\text{amine: }\mathrm{C{-}N}
\end{equation}
Look for the carbonyl first. A nitrogen next to a carbonyl is an amide, not an amine.

\subsection{Exam recipe: functional groups in a molecule}
\textbf{Use when:} an organic structure is shown and the functional groups present must be identified.

\textbf{Given / Find:} Given a line, condensed, or displayed structure; find every functional group named in the question or answer options.

\textbf{Required references:} Table~\ref{tab:organic:functional-groups} and Equation~\eqref{eq:organic:heteroatom-distinction}; use Equation~\eqref{eq:organic:aromatic-motif} for aromatic rings.

\textbf{Steps:}
\begin{enumerate}
    \item Mark every heteroatom \(\ce{O}\), \(\ce{N}\), halogen, triple bond, and carbonyl \(\ce{C=O}\).
    \item For each carbonyl carbon, inspect its directly attached atom and classify acid, ester, amide, aldehyde, or ketone.
    \item Classify remaining \(\ce{O}\) and \(\ce{N}\) motifs as alcohol, ether, or amine only after carbonyl groups are assigned.
    \item Identify aromatic rings and nitriles separately.
    \item Match the complete list against the requested response or eliminate options containing a missing motif.
\end{enumerate}

\textbf{Checks:} Every reported group must correspond to an explicit motif; an acid counts as carboxylic acid, not separate alcohol plus ketone.

\textbf{Common traps:} Calling an ester an ether; calling an amide an amine; overlooking a carbonyl or an aromatic ring.

\subsection{Isomerism, chirality, and biomolecules}

\subsubsection{Structural isomers}
\begin{equation}
\label{eq:organic:structural-isomers}
\text{structural isomers}=\text{same molecular formula, different connectivity}
\end{equation}
Structural isomers can have different names, boiling points, and functional groups.

\subsection{Exam recipe: molecular formula consistency and constitutional isomer enumeration}
\textbf{Use when:} a displayed structure must be tested against a molecular formula, or the number of constitutional isomers with a stated functional group is requested.

\textbf{Given / Find:} Given a molecular formula, candidate structures, or a functional-group restriction such as alcohol; find inconsistent structures or the number of distinct connectivities.

\textbf{Required references:} Equations~\eqref{eq:organic:dbe}, \eqref{eq:organic:structural-isomers}, and Table~\ref{tab:organic:functional-groups}.

\textbf{Steps:}
\begin{enumerate}
    \item Calculate DBE from Equation~\eqref{eq:organic:dbe}; use it to determine how many total rings and multiple bonds a valid candidate can contain.
    \item To test a shown option, count each element in the drawing and compare with the stated formula, then reject it if the atom count or DBE does not match.
    \item To enumerate isomers, first list distinct carbon skeletons compatible with the formula and DBE, without placing the required functional group twice on symmetry-equivalent positions.
    \item Place the required functional group on each nonequivalent permitted position of each skeleton; for an alcohol, attach \(\ce{-OH}\) to carbon and exclude ethers unless the question includes them.
    \item Count only structures with different connectivity as required by Equation~\eqref{eq:organic:structural-isomers}; do not count mirror images or rotations as additional constitutional isomers.
\end{enumerate}

\textbf{Checks:} Every counted structure has the stated formula, DBE, and required functional group; for saturated acyclic \(\ce{C4H10O}\) alcohols, DBE must be zero.

\textbf{Common traps:} counting an ether among alcohols; counting the same structure redrawn from another end; forgetting a branched carbon skeleton.

\subsubsection{Chiral center}
\begin{equation}
\label{eq:organic:chiral-center}
\mathrm{C^*}=\text{tetrahedral carbon attached to four different groups}
\end{equation}
A molecule with one chiral center is chiral. Symmetry can make molecules with multiple stereocenters achiral.

\subsubsection{Amino acids and peptide bond}
\begin{equation}
\label{eq:organic:amino-acid}
\text{amino acid motif: }\mathrm{H_2N{-}CHR{-}COOH}
\end{equation}
\begin{equation}
\label{eq:organic:peptide-bond}
\text{peptide bond: }\mathrm{C(=O){-}NH}
\end{equation}
Amino acids contain amine and carboxylic acid groups. Peptides are connected by amide bonds.

\subsection{Exam recipe: chirality and peptide or biomolecule recognition}
\textbf{Use when:} a structure must be classified as chiral, amino-acid-like, or peptide-containing.

\textbf{Given / Find:} Given a molecular structure; find stereogenic carbons and identify amino acid or peptide motifs.

\textbf{Required references:} Equations~\eqref{eq:organic:chiral-center}, \eqref{eq:organic:amino-acid}, and \eqref{eq:organic:peptide-bond}; Table~\ref{tab:organic:functional-groups}.

\textbf{Steps:}
\begin{enumerate}
    \item Locate tetrahedral carbon candidates and list the four attached groups for each candidate.
    \item Mark \(\mathrm{C^*}\) only when all four groups differ; if multiple centers occur, check for internal symmetry before concluding that the molecule is chiral.
    \item Search for the amino acid pattern: amine and carboxylic acid attached through the same alpha-carbon framework.
    \item Search for \(\mathrm{C(=O){-}NH}\) links between residues; classify each such link as a peptide (amide) bond.
    \item State only the requested property: number of chiral centers, chirality, amino acid motif, or peptide bond count.
\end{enumerate}

\textbf{Checks:} A carbon with two identical substituents is not stereogenic; a peptide bond must include a carbonyl directly bonded to nitrogen.

\textbf{Common traps:} Counting a planar carbonyl carbon as chiral; calling any \(\ce{C-N}\) bond a peptide bond; treating an amide nitrogen as an amine group.

\subsection{Polymers}

\subsubsection{Degree of polymerization}
\begin{equation}
\label{eq:organic:polymerization-degree}
n=\frac{M_{\mathrm{polymer}}}{M_{\mathrm{repeat}}}
\end{equation}
The degree of polymerization is the number of repeat units in the chain.

\subsection{Exam recipe: polymer repeat unit and degree of polymerization}
\textbf{Use when:} a polymer structure, monomer, repeat-unit mass, or polymer molar mass is given.

\textbf{Given / Find:} Given the repeat-unit structure or its molar mass and \(M_{\mathrm{polymer}}\); find the repeat unit, \(M_{\mathrm{repeat}}\), or degree \(n\).

\textbf{Required references:} Equation~\eqref{eq:organic:polymerization-degree} and the atomic molar masses supplied in the problem or data sheet.

\textbf{Steps:}
\begin{enumerate}
    \item Identify the smallest structural segment that repeats along the chain and write its atom count.
    \item Compute \(M_{\mathrm{repeat}}\) from that segment unless it is supplied.
    \item Insert consistent molar-mass units into Equation~\eqref{eq:organic:polymerization-degree}.
    \item Calculate \(n\) and round only as justified by the idealized chain or stated precision.
\end{enumerate}

\textbf{Checks:} \(n\) is dimensionless and is normally close to a positive integer in an idealized exam problem; multiplying \(nM_{\mathrm{repeat}}\) should recover \(M_{\mathrm{polymer}}\).

\textbf{Common traps:} Using monomer mass when condensation removes a small molecule; counting end groups as part of every repeat unit; reporting \(n\) with mass units.
